\section*{第 \ref{ch:g10:concentration-and-dilution} 章——浓度与稀释}

\begin{solution}{exo:g10:concentration-and-dilution:1}
$C_m = \qty{5.0}{g} / \qty{0.250}{L} = \qty{20}{g/L}$.
\end{solution}

\begin{solution}{exo:g10:concentration-and-dilution:2}
$c = \qty{0.050}{mol} / \qty{0.200}{L} = \qty{0.25}{mol/L}$.
\end{solution}

\begin{solution}{exo:g10:concentration-and-dilution:3}
$M(\ce{CuSO4}) = 63.5 + 32.1 + 4 \times 16.0 = \qty{159.6}{g/mol}$;
$n = 0.10 \times 0.1000 = \qty{0.0100}{mol}$; $m = 0.0100 \times 159.6 =
\qty{1.60}{g}$.
\end{solution}

\begin{solution}{exo:g10:concentration-and-dilution:4}
$F = 1.0 / 0.050 = 20$; $V_0 = \qty{100.0}{mL} / 20 = \qty{5.0}{mL}$.
\end{solution}

\begin{solution}{exo:g10:concentration-and-dilution:5}
$C_m = c \times M = 0.10 \times 58.5 = \qty{5.85}{g/L}$，约
\qty{5.9}{g/L}.
\end{solution}

\begin{solution}{exo:g10:concentration-and-dilution:6}
$V_0 = c_1 V_1 / c_0 = 0.020 \times 250.0 / 0.50 = \qty{10.0}{mL}$
（稀释倍数为 25）。需用 \qty{10.0}{mL} 移液管和
\qty{250.0}{mL} 容量瓶。
\end{solution}

\begin{solution}{exo:g10:concentration-and-dilution:7}
容量瓶只有一条标线，只为一个体积而制，
误差远小于百分之一；烧杯的刻度
只是粗略的参考。同样，移液管量取一个体积
非常精确，而量筒较粗，读数
精度较低。
\end{solution}

\begin{solution}{exo:g10:concentration-and-dilution:8}
介于 \qty{0.02}{mol/L} 与 \qty{0.04}{mol/L} 之间：约 \qty{0.03}{mol/L}。
若要更精确，可在两者之间再配制标准溶液，例如在
0.025、0.030 和 \qty{0.035}{mol/L} 处配制，在需要处设置更细的标准色阶。
\end{solution}

\begin{solution}{exo:g10:concentration-and-dilution:9}
$M(\ce{C6H12O6}) = \qty{180.0}{g/mol}$; $c = 50 / 180.0 =
\qty{0.28}{mol/L}$.
\end{solution}

\begin{solution}{exo:g10:concentration-and-dilution:10}
$M = \qty{342.0}{g/mol}$，故 $n = \qty{1.00}{mol}$，$c =
\qty{1.00}{mol/L}$。稀释十倍后：$c = \qty{0.100}{mol/L}$，即
$C_m = 0.100 \times 342.0 = \qty{34.2}{g/L}$。一小勺
\qty{0.020}{L} 的溶液含蔗糖 $34.2 \times 0.020 = \qty{0.68}{g}$。
\end{solution}

\begin{solution}{exo:g10:concentration-and-dilution:11}
(a) 溶质质量相同，体积为两倍：\qty{6.0}{g/L}。
(b) 体积约减半，溶质质量不变：约
\qty{24}{g/L}.
\end{solution}

\begin{solution}{exo:g10:concentration-and-dilution:12}
\begin{enumerate}
  \item 三次十倍稀释共稀释 $10^3$ 倍：
        $\qty{0.10}{mol/L} / 1000 = \qty{1.0e-4}{mol/L}$.
  \item 若一次稀释 1000 倍，则要么需量取极少量的母液
        （\qty{0.1}{mL} 加入 \qty{100}{mL} 中），难以用移液管准确量取，
        要么需用极大的容量瓶（\qty{1}{mL} 加入
        \qty{1}{L} 中）。三次十倍稀释只需普通的移液管
        和容量瓶。
\end{enumerate}
\end{solution}

\begin{solution}{exo:g10:concentration-and-dilution:13}
\begin{enumerate}
  \item $M(\ce{CaCl2}) = 40.1 + 2 \times 35.5 = \qty{111.1}{g/mol}$;
        $n = 11.1 / 111.1 = \qty{0.100}{mol}$; $c = 0.100 / 0.5000 =
        \qty{0.200}{mol/L}$.
  \item 每个 \ce{CaCl2} 电离出一个 \ce{Ca^2+} 和两个 \ce{Cl-}：
        $[\ce{Ca^2+}] = \qty{0.200}{mol/L}$，且 $[\ce{Cl-}] =
        \qty{0.400}{mol/L}$.
\end{enumerate}
\end{solution}

\begin{solution}{exo:g10:concentration-and-dilution:14}
\begin{enumerate}
  \item $r = \qty{0.60}{cm}$, $h = \qty{0.20}{cm}$:
        $V = \pi \times 0.60^2 \times 0.20 = \qty{0.23}{cm^3}$，即多出
        \qty{0.23}{mL} 的水。
  \item 溶质同样溶于 \qty{100.23}{mL} 中，而非 \qty{100.00}{mL} 中：
        浓度偏低 $0.23 / 100.23 \approx
        \qty{0.23}{\percent}$.
\end{enumerate}
\end{solution}

\begin{solution}{exo:g10:concentration-and-dilution:15}
$n = 0.20 \times 0.100 + 0.10 \times 0.300 = 0.020 + 0.030 =
\qty{0.050}{mol}$，溶于 \qty{0.400}{L} 中：$c = \qty{0.125}{mol/L}$。这不是
并非简单平均 \qty{0.15}{mol/L}：稀溶液的体积是浓溶液的三倍，
因此混合物更接近 \qty{0.10}{mol/L}。这是
以体积为权重的平均值。
\end{solution}

\begin{solution}{pb:g10:concentration-and-dilution:1}
\textbf{1.} \qty{0.9}{g} 溶于 \qty{0.100}{L}：$C_m = \qty{9.0}{g/L}$。

\textbf{2.} \qty{9.0}{g} 溶于 \qty{1000}{mL}，即每
毫升含 \qty{9.0}{mg}。

\textbf{3.} $m = 9.0 \times 0.500 = \qty{4.5}{g}$.

\textbf{4.} 是的：溶液中每升约含 \qty{9}{g}，饱和时每千克水可溶解
\qty{360}{g}，溶液浓度约低四十倍。

\textbf{5.} 共需盐 $9.0 \times 1.0 = \qty{9.0}{g}$。

\textbf{6.} $M(\ce{NaCl}) = 23.0 + 35.5 = \qty{58.5}{g/mol}$.

\textbf{7.} $c = 9.0 / 58.5 = \qty{0.154}{mol/L}$.

\textbf{8.} $n = 0.154 \times 0.500 = \qty{0.0769}{mol}$（或
$4.5 / 58.5$）。

\textbf{9.} 每个 \ce{NaCl} 电离出一个 \ce{Na+} 和一个 \ce{Cl-}，二者浓度均为
\qty{0.154}{mol/L}。

\textbf{10.} $0.0769 \times \num{6.02e23} = \num{4.63e22}$ 个钠离子。

\textbf{11.} $F = 200 / 9.0 = 22.2$.

\textbf{12.} 袋中的 \qty{4.5}{g} 全部取自浓
溶液：$V_0 = \qty{4.5}{g} / \qty{200}{g/L} = \qty{0.0225}{L} =
\qty{22.5}{mL}$。

\textbf{13.} $V_0 = V_1 / F = 500 / 22.2 = \qty{22.5}{mL}$.

\textbf{14.} $c_0 = 200 / 58.5 = \qty{3.42}{mol/L}$;
$c_0 V_0 = 3.42 \times 0.0225 = \qty{0.0769}{mol}$ 与
$c_1 V_1 = 0.154 \times 0.500 = \qty{0.0769}{mol}$：相等。

\textbf{15.} 约加无菌水 $500 - 22.5 = \qty{478}{mL}$。

\textbf{16.} 用 \qty{25}{mL} 的刻度移液管（或滴定管），
读数精确至 \qty{0.1}{mL}。

\textbf{17.} 用 \qty{500.0}{mL} 容量瓶配制，定容至刻度线。

\textbf{18.} $C_m = 200 \times 23.0 / 500 = \qty{9.20}{g/L}$，偏浓
约 $0.20 / 9.0 \approx \qty{2.2}{\percent}$。

\textbf{19.} 取 \qty{22.5}{mL} 的浓（“20 \%”）溶液，
用于配制一袋 \qty{500}{mL} 的溶液。
\end{solution}
